\section{Finite Turbo SPIN Certificate}
\label{app:turbo-proof}

This appendix specifies the fixed maps and the first-moment interfaces
behind Theorem~\ref{thm:turbo-spin}. The bounds average over independently
sampled outer column maps, packet and regional permutations, and state
updates. They concern the ideal setup distribution, rather than the
distance of a particular seeded implementation.

\paragraph{Outer code and expected measure.}
Represent $\F_{16}$ in the polynomial basis with modulus
$X^4+X+1$. Each RS row interpolates its four input values at
$0,1,2,3\in\F_{16}$ and evaluates at $0,\ldots,7$.
Write $Q=2^{32}$. For a group, the number $A_h$ of nonzero outer words
with exactly $h$ nonzero $32$-bit columns is
\begin{equation}
 A_h=\binom8h\sum_{j=0}^{h-5}(-1)^j\binom hj
       \bigl(Q^{h-4-j}-1\bigr),\qquad 5\le h\le8.
 \label{eq:turbo-column-spectrum}
\end{equation}
Indeed, shortening the eight RS rows to a specified set of $h$ columns
leaves $8(h-4)$ dimensions over $\F_{16}$; inclusion--exclusion requires
all selected columns to be nonzero.

For each nonzero column, multiplication by a uniform nonzero element of
$\F_{2^{32}}$ sends any fixed nonzero vector uniformly to the nonzero
$32$-bit vectors. The adjoint family has the same property: if two
coefficients send one fixed nonzero vector to the same image, their
difference gives a nonzero field multiplication with zero adjoint image,
which is impossible. Let $C_v$ be the expected number of nonzero group
messages yielding $v$ active four-bit packets before the group shuffle.
Then
\begin{equation}
  P(z):=(1+15z)^8-1,\qquad
  \sum_{v=5}^{64}C_vz^v
   =\sum_{h=5}^{8}A_h\left(\frac{P(z)}{Q-1}\right)^h.
  \label{eq:turbo-packet-spectrum}
\end{equation}
The coefficient of $z^v$ is an expected message count, not a probability
or the realized spectrum of a sampled group. The sum of the coefficients
is $2^{128}-1$. Given $v$, the group shuffle makes the packet support
uniform among $\binom{64}{v}$ choices, with uniform nonzero labels on the
active packets.

\paragraph{Fixed inner maps.}
For $s=15$, index the $64$ output positions by
$z=(z_0,\ldots,z_5)\in\F_2^6$, with coordinate number
$\sum_i2^iz_i$. The columns of $A$ are the truth tables of
$1,z_0,\ldots,z_5$ and the following eight functions, in order:
\[
\begin{gathered}
 z_0z_1,\quad z_0z_3,\quad z_4z_5,\quad
 z_0z_5+z_2z_3,\quad z_1z_3+z_2z_4,\\
 z_0z_4+z_2z_5,\quad z_1z_5+z_3z_4,\quad
 z_1z_2+z_3z_5.
\end{gathered}
\]
This is the paired $15$-state map. For $s=22$, use the map $A$
specified in Appendix~\ref{app:shared-inner-map}.
In both cases the feedback is $A^\top$. Each step uses an independent
uniform invertible state map. These are separate certified maps, not
instances of the paper's BCH--IMT transvection family.

\paragraph{First moment and occupation coverage.}
Let $Z_H$ count nonzero messages with output weight at most $H$, and let
$q$ count their nonzero groups. Write $W_v$ for an upper bound on the
conditional output-weight event for one group with $v$ active packets,
averaged over its support, labels, route, and inner updates. Likewise,
$V_{u,v}$ bounds the corresponding two-group event. The first two
contributions obey
\begin{equation}
 U_1=L\sum_{v=5}^{64}C_vW_v,\qquad
 U_2=\binom L2\sum_{5\le u\le v\le64}
       (2-\1_{u=v})C_uC_vV_{u,v}.
 \label{eq:turbo-sparse-fold}
\end{equation}
The factors count unordered group positions and both support assignments
when $u\ne v$. The two-group computation keeps the common route and the
inner state across every region.

For larger occupations, the numerical calculation compares Turbo SPIN
with a reference outer that maps each $128$-bit group through four
RS$[16,8]$ rows over $\F_{16}$. Its $16$ aligned $16$-bit columns receive
independent nonzero $\F_{2^{16}}$ scalar randomizers. It then uses the
same $64$ four-bit packet positions, route, and fixed inner map as Turbo
SPIN at the selected length. Complete numerical bounds for this reference
outer were checked for every occupation; the following pointwise
comparison transfers them without assuming identical outer spectra.

For a fixed nonzero $256$-bit vector with $h$ nonzero columns, the
number of RS words with that column support is at most $(Q-1)^{h-4}$.
Independent column randomization therefore bounds its expected atom count
by $(Q-1)^{-4}$. Relative to uniform input, the outer-group measure has
pointwise multiplier
\[
 \beta_{\rm Turbo}=\frac{2^{256}}{(2^{32}-1)^4}
 <\frac{2^{256}}{(2^{16}-1)^8}
 =\beta_{\rm RS16}.
\]
For each fixed active-group set, this domination transfers the retained
uniform-input and routing-conditioned bounds for $q\ge3$, with the same
inner and route. The transfer is over the complete outer measure;
one-column marginals alone would not justify it.

The old RS$[16,8]$ outer has no words with fewer than nine active packets,
whereas Turbo SPIN can have five through eight. This difference does not
matter: the one- and two-group sums \eqref{eq:turbo-sparse-fold} are
computed for Turbo SPIN itself and include those supports at every length.

\paragraph{Bounding one large occupation.}
For $q\ge3$, the regional bound needs the coefficient of $x^q$ in a
product of nonnegative two-step matrix polynomials, normalized by
$\binom Lq$. When $q$ is small, we compute that coefficient exactly by
repeated squaring, truncated at degree $q$. When $q$ is large, the exact
coefficient is expensive, and we use a positive marker $t$. Every
coefficient of a polynomial $P$ with nonnegative entries satisfies
$[x^q]P(x)\le P(t)/t^q$, so one matrix power at $x=t$ replaces the
coefficient. The routing-conditioned variant first raises each two-step
operator entrywise to a power $\alpha\le1$; the witnesses use
$\alpha\in\{1,4/5,7/10,3/5\}$ and switch from exact coefficients to
markers above $q=279$ at $K=2^{18}$ and above $q=315$ at $K=2^{20}$.

\paragraph{Margins.}
The occupation sets are $q=1,\ldots,L$ with $L=512,2048,8192$ for the
three lengths. Summing all occupations gives failure bounds of
$2^{-61.884}$, $2^{-58.2408}$, and $2^{-56.2570}$. At the two larger
lengths the one-group term dominates: it alone is about $2^{-58.2456}$
and $2^{-56.2570}$. The exact cutoff weights are $13107$, $52428$, and
$209715$; hence the integer minimum-distance lower bounds outside the
failure event are $13108$, $52429$, and $209716$.

\paragraph{Formal verification.}
For $K=2^{18}$ and $K=2^{20}$, a Lean~4 development is the authority for
Theorem~\ref{thm:turbo-spin}\repositorycite. It models the ideal setup
distribution of Section~\ref{sec:turbo-spin}, with $s=22$, for every
injective choice of the eight RS evaluation points. It proves
\[
 \Prob_\theta\bigl[\exists x\ne0:\wt(E_\theta(x))\le H\bigr]
 \le 2^{-58.24083}
 \quad\text{and}\quad
 \le 2^{-56.25709},
\]
respectively. Lean checks every numerical certificate behind these bounds,
for all occupations, including the integer rounding and scaling that this
appendix omits. The proofs use only Lean's standard axioms: propositional
extensionality, quotient soundness, and choice. The rate statement follows
from Lemma~\ref{lem:structured-inner-bijective} and is outside the formal
statement. The $K=2^{16}$ bound is not formalized; it rests on the
outward-arithmetic replay recorded in the repository.
